\documentclass[conference]{IEEEtran}

\usepackage{amsmath,amssymb}
\usepackage{graphicx}
\usepackage{cite}
\usepackage{url}
\usepackage{booktabs}
\usepackage{algorithm}
\usepackage{algorithmic}
\usepackage{listings}
\usepackage{xcolor}
\usepackage[hidelinks,colorlinks=true,urlcolor=blue]{hyperref}

\lstset{
  basicstyle=\ttfamily\scriptsize,
  keywordstyle=\color{blue},
  commentstyle=\color{gray},
  language=C++,
  numbers=left,
  numberstyle=\tiny\color{gray},
  numbersep=4pt,
  frame=single,
  breaklines=true,
  tabsize=2,
  captionpos=b
}

\begin{document}

\title{OpenMP Parallelisation of a Two-Dimensional\\
Ising Model Monte Carlo Simulation}

\author{
  \IEEEauthorblockN{Anushka Chauhan}
  \IEEEauthorblockA{
    Indian Institute of Technology Mandi\\
    Roll No: B23310\\
    \href{https://github.com/03anushka03/HPSC\_Project.git}{GitHub Repository}
  }
}

\maketitle

\begin{abstract}
A two-dimensional Ising model is implemented and parallelised using OpenMP.
The solver uses the Metropolis Monte Carlo algorithm on a square lattice with
periodic boundary conditions. Physical observables including magnetisation,
energy, heat capacity, and susceptibility are computed over a range of
temperatures. The implementation is verified against the exact analytical
critical temperature $T_c \approx 2.269\,J/k_B$. Two parallelisation
strategies are explored: an embarrassingly parallel temperature sweep and a
checkerboard lattice decomposition. Parallelisation improves performance for
larger workloads, with speedup increasing with the number of temperature
points being simulated.
\end{abstract}

\begin{IEEEkeywords}
Ising model, Monte Carlo, Metropolis algorithm, OpenMP, phase transition,
parallel computing, checkerboard decomposition.
\end{IEEEkeywords}

% ─────────────────────────────────────────────────────────────
\section{Introduction}

The Ising model is one of the most studied models in statistical physics.
It describes a system of interacting spins on a lattice and exhibits a
phase transition between ordered and disordered states. Despite its
simplicity, it captures the essential physics of ferromagnetism and
serves as a standard benchmark for numerical methods.

From a computational standpoint, Monte Carlo simulation of the Ising model
is expensive. A single sweep over the lattice involves $N = L^2$ random
spin flip attempts, each requiring evaluation of the local energy change.
Running thousands of such sweeps over many temperature values leads to
significant runtime.

OpenMP is applied in this work to parallelise the simulation. The goals are
to reduce runtime while maintaining physical correctness, and to analyse
scaling behaviour with thread count.

% ─────────────────────────────────────────────────────────────
\section{Mathematical Model}

Each site $(i,j)$ on an $L \times L$ lattice carries a spin
$s_{i,j} \in \{+1, -1\}$. The Hamiltonian is

\begin{equation}
  H = -J \sum_{\langle i,j \rangle} s_i s_j,
  \label{eq:H}
\end{equation}

where $J > 0$ is the coupling constant and the sum runs over nearest-neighbour
pairs. Periodic boundary conditions are applied so that
$s_{L,j} \equiv s_{0,j}$ and $s_{i,L} \equiv s_{i,0}$.

\subsection{Contact Force Analogy}

When two spins interact, the local energy change due to flipping spin
$s_{i,j}$ is

\begin{equation}
  \Delta E = 2J\, s_{i,j} \sum_{\mathrm{nn}} s_k,
  \label{eq:dE}
\end{equation}

where the sum is over the four nearest neighbours.
Since $\Delta E \in \{-8J, -4J, 0, 4J, 8J\}$, only five distinct
Boltzmann factors need to be evaluated per temperature, so they are
pre-computed before each sweep.

\subsection{Thermodynamic Observables}

The magnetisation per spin and energy per spin are

\begin{equation}
  \langle|M|\rangle = \frac{1}{N}\left\langle\left|\sum_{i,j}s_{i,j}\right|\right\rangle,
  \qquad
  \langle E/N \rangle = \frac{\langle H \rangle}{N}.
\end{equation}

Heat capacity and susceptibility follow from fluctuations:

\begin{equation}
  C_v = \frac{\langle E^2\rangle - \langle E\rangle^2}{k_B T^2},
  \qquad
  \chi = \frac{N\!\left(\langle M^2\rangle - \langle M\rangle^2\right)}{k_B T}.
  \label{eq:cv_chi}
\end{equation}

The exact Onsager critical temperature \cite{onsager1944} is

\begin{equation}
  T_c = \frac{2J}{k_B\ln(1+\sqrt{2})} \approx 2.269\,\frac{J}{k_B}.
  \label{eq:tc}
\end{equation}

% ─────────────────────────────────────────────────────────────
\section{Numerical Algorithm}

The Metropolis Monte Carlo algorithm \cite{metropolis1953} generates a Markov
chain whose stationary distribution is the Boltzmann distribution. One sweep
consists of $N = L^2$ proposed single-spin flips:

\begin{algorithm}
\caption{Metropolis sweep}
\begin{algorithmic}[1]
\FOR{step $= 1$ \TO $N$}
  \STATE Select site $(i,j)$ uniformly at random
  \STATE Compute $\Delta E$ from \eqref{eq:dE}
  \IF{$\Delta E < 0$ \OR $r < e^{-\Delta E/T}$, $r\sim\mathcal{U}(0,1)$}
    \STATE Flip: $s_{i,j} \leftarrow -s_{i,j}$
  \ENDIF
\ENDFOR
\end{algorithmic}
\end{algorithm}

Each temperature point runs 5\,000 equilibration sweeps followed by
10\,000 measurement sweeps. Observables are accumulated during the
measurement phase and averaged at the end.

The standard \texttt{rand()} function is not used because it is not
thread-safe. Each thread instead owns a private \texttt{std::mt19937}
Mersenne Twister instance seeded from a master generator.

% ─────────────────────────────────────────────────────────────
\section{Verification}

Three tests are used to verify the implementation.

\textit{Low-temperature behaviour.} At $T = 1.0\,J/k_B$ the system should
be almost fully ordered. The simulation gives $|\langle M\rangle| = 0.9993$
and $E/N = -1.997$, both within 0.3\% of the expected ground state values
of 1 and $-2J$ respectively.

\textit{High-temperature behaviour.} At $T = 5.0\,J/k_B$ the system should
be disordered. The simulation gives $|\langle M\rangle| = 0.028$, consistent
with a paramagnetic phase.

\textit{Critical temperature.} The magnetisation drops sharply and $\chi$
peaks near $T \approx 2.3\,J/k_B$, in good agreement with $T_c = 2.269$.
The small offset is a known finite-size effect for $L = 50$
\cite{ferdinand1969}.

\begin{table}[h]
\centering
\caption{Verification summary ($L=50$, $J=k_B=1$)}
\label{tab:verify}
\begin{tabular}{lcc}
\toprule
Test & Expected & Simulated \\
\midrule
$|\langle M\rangle|$ at $T=1.0$ & $\approx 1.0$ & $0.9993$ \\
$E/N$ at $T=1.0$ & $\approx -2.0$ & $-1.997$ \\
$|\langle M\rangle|$ at $T=5.0$ & $\approx 0$ & $0.028$ \\
$T_c$ (peak of $\chi$) & $2.269$ & $\approx 2.3$ \\
\bottomrule
\end{tabular}
\end{table}

Fig.~\ref{fig:mag} shows the magnetisation transition and
Fig.~\ref{fig:energy} shows the energy curve. Both match the expected
qualitative behaviour. The convergence of the statistical error with
measurement sweeps is shown in Fig.~\ref{fig:conv}, confirming
$\mathcal{O}(1/\sqrt{N})$ behaviour as expected for Monte Carlo.

\begin{figure}[t]
  \centering
  \includegraphics[width=\columnwidth]{fig_mag.pdf}
  \caption{Magnetisation vs temperature. The sharp drop near
           $T_c = 2.269\,J/k_B$ (dashed) marks the phase transition.}
  \label{fig:mag}
\end{figure}

\begin{figure}[t]
  \centering
  \includegraphics[width=\columnwidth]{fig_energy.pdf}
  \caption{Energy per spin vs temperature. The smooth S-shaped curve is
           characteristic of the second-order phase transition.}
  \label{fig:energy}
\end{figure}

\begin{figure}[t]
  \centering
  \includegraphics[width=\linewidth]{fig_cv_chi.pdf}
  \caption{Heat capacity $C_v$ (left) and susceptibility $\chi$ (right).
           Both peak near $T_c$, with the offset from the exact value
           being a finite-size effect.}
  \label{fig:cvchi}
\end{figure}

\begin{figure}[t]
  \centering
  \includegraphics[width=\linewidth]{fig_snapshots.pdf}
  \caption{Spin configurations on the $50\times50$ lattice. Blue $=-1$,
           red $=+1$. Ordered phase (left), critical fluctuations (centre),
           and disordered phase (right).}
  \label{fig:snaps}
\end{figure}

\begin{figure}[t]
  \centering
  \includegraphics[width=\columnwidth]{fig_convergence.pdf}
  \caption{Statistical error in $|\langle M\rangle|$ vs number of
           measurement sweeps. The error follows $\mathcal{O}(1/\sqrt{N})$
           as expected for Monte Carlo.}
  \label{fig:conv}
\end{figure}

% ─────────────────────────────────────────────────────────────
\section{Profiling}

Profiling was carried out to identify which parts of the code take the most
time. It was found that the Metropolis sweep dominates the total runtime.

This is expected: for each of the 41 temperature points, the inner loop
performs $N = 2{,}500$ spin-flip attempts per sweep, repeated
15\,000 times (5\,000 equilibration + 10\,000 measurement). The total
number of spin-flip evaluations is approximately $1.5 \times 10^9$.

The observable calculation (\texttt{calculate\_magnetization} and
\texttt{calculate\_energy}), file I/O, and lattice initialisation together
account for less than 7\% of runtime. These are not the primary target
for parallelisation.

The total serial runtime for 41 temperature points on an $L=50$ lattice
was measured at $\mathbf{61{,}274}$~ms.

% ─────────────────────────────────────────────────────────────
\section{Parallelisation}

OpenMP is used to parallelise the simulation. Two strategies are implemented.

\subsection{Strategy 1: Parallel Temperature Loop}

Each temperature point is an independent simulation with its own copy of the
lattice. These can be run in parallel without any shared state:

\begin{lstlisting}[caption={Strategy 1 — parallel temperature sweep}]
#pragma omp parallel for schedule(dynamic)
for (int ti = 0; ti < n_temps; ++ti) {
    double T = T_min + ti * dT;
    vector<vector<int>> spin(L, vector<int>(L));
    initialize_lattice(spin, L, tid);
    run_simulation(spin, L, T, J, ...);
}
\end{lstlisting}

Dynamic scheduling is used because the equilibration cost varies with
temperature (critical slowing-down is more severe near $T_c$).

\subsection{Strategy 2: Checkerboard Decomposition}

For a single large lattice, spins cannot all be updated simultaneously
because updating one spin changes the energy of its neighbours, leading to
race conditions. The checkerboard (red-black) scheme avoids this by
partitioning the lattice into two sub-lattices where no two sites of the
same colour are neighbours of each other:

\begin{lstlisting}[caption={Strategy 2 — checkerboard sweep}]
for (int colour = 0; colour < 2; ++colour) {
    #pragma omp parallel for schedule(static)
    for (int i = 0; i < L; ++i) {
        for (int j = (i+colour)%2; j < L; j+=2) {
            // all same-colour sites are independent
            double dE = compute_deltaE(spin,i,j,L,J);
            if (dE<0 || rand01(tid)<exp(-dE/T))
                spin[i][j] *= -1;
        }
    }
} // implicit barrier between colour passes
\end{lstlisting}

The implicit OpenMP barrier at the end of each \texttt{parallel for}
ensures correctness: the red sub-lattice is fully updated before the
black sub-lattice begins, preserving detailed balance.

Thread safety is achieved by assigning each thread its own
\texttt{std::mt19937} random number generator. The original code used a
global \texttt{rand()} call and a global timing accumulator
(\texttt{metro\_time}); both are race conditions and have been removed.
File writes are protected by \texttt{\#pragma omp critical}.

% ─────────────────────────────────────────────────────────────
\section{Performance Analysis}

Speedup and efficiency are defined as

\begin{equation}
  S(p) = \frac{T_1}{T_p}, \qquad E(p) = \frac{S(p)}{p},
\end{equation}

where $T_p$ is wall-clock time with $p$ threads.

Profiling showed that approximately 93\% of runtime is in the Metropolis
sweep, giving a parallel fraction $f = 0.93$. Amdahl's Law then predicts a
theoretical maximum speedup of $1/(1-f) \approx 14\times$.

Fig.~\ref{fig:speedup} shows the strong-scaling behaviour and
Fig.~\ref{fig:runtime} shows the wall-clock times.
For small thread counts the speedup is close to ideal. Efficiency drops
at higher thread counts due to thread management overhead and the serial
fraction in observable computation and I/O.

\begin{figure}[t]
  \centering
  \includegraphics[width=\linewidth]{fig_speedup.pdf}
  \caption{Speedup $S(p)$ (left) and efficiency $E(p)$ (right) vs thread
           count for Strategy~1. Dashed lines show ideal linear scaling.}
  \label{fig:speedup}
\end{figure}

\begin{figure}[t]
  \centering
  \includegraphics[width=\columnwidth]{fig_runtime.pdf}
  \caption{Wall-clock time for serial and parallel ($p=2$) runs at
           different workload sizes. Parallel provides clear benefit at
           larger numbers of temperature points.}
  \label{fig:runtime}
\end{figure}

For small workloads (a single temperature point) the parallel version is
slightly slower than serial due to thread creation and synchronisation
overhead. This is the same behaviour observed in the DEM solver
\cite{dem_report}: parallelisation only pays off when there is enough
work per thread.

% ─────────────────────────────────────────────────────────────
\section{Conclusion}

The Ising model Monte Carlo solver produces correct results for all
verification tests. The magnetisation, energy, heat capacity, and
susceptibility all behave as expected, with the phase transition observed
near the exact Onsager value $T_c = 2.269\,J/k_B$.

Parallelisation using OpenMP improves performance, especially when many
temperature points are simulated in parallel. The embarrassingly parallel
temperature sweep (Strategy~1) is more effective for this workload than
the checkerboard decomposition, since it requires no synchronisation inside
the inner loop. However, Strategy~2 remains useful when a single large
lattice must be updated and memory constraints prevent replication.

Further improvements could include cluster algorithms (Wolff
\cite{wolff1989}) to eliminate critical slowing-down, or reducing the
observable calculation serial fraction using \texttt{omp reduction}.

\section*{Acknowledgements}

Some help from large language models (ChatGPT and Claude) was used for
debugging, plot generation, and \LaTeX{} fixes.

% ─────────────────────────────────────────────────────────────
\begin{thebibliography}{9}

\bibitem{onsager1944}
L.~Onsager,
``Crystal statistics. I. A two-dimensional model with an order-disorder
transition,''
\emph{Phys.\ Rev.}, vol.~65, pp.~117--149, 1944.

\bibitem{metropolis1953}
N.~Metropolis \emph{et al.},
``Equation of state calculations by fast computing machines,''
\emph{J.\ Chem.\ Phys.}, vol.~21, no.~6, pp.~1087--1092, 1953.

\bibitem{ferdinand1969}
A.~E.~Ferdinand and M.~E.~Fisher,
``Bounded and inhomogeneous Ising models,''
\emph{Phys.\ Rev.}, vol.~185, pp.~832--846, 1969.

\bibitem{wolff1989}
U.~Wolff,
``Collective Monte Carlo updating for spin systems,''
\emph{Phys.\ Rev.\ Lett.}, vol.~62, pp.~361--364, 1989.

\bibitem{dem_report}
A.~Chauhan,
``OpenMP parallelisation of a three-dimensional discrete element method
solver,''
IIT Mandi, B23310, 2025.

\bibitem{newman1999}
M.~E.~J.~Newman and G.~T.~Barkema,
\emph{Monte Carlo Methods in Statistical Physics}.
Oxford: Clarendon Press, 1999.

\end{thebibliography}

\end{document}
